\section{Related work}\label{sec:related}

Our literature search was targeted, not exhaustive. The comparisons below refer to the statements we checked in the cited sources.

\subsection{Spanning trees and determinantal distributions in probabilistic circuits}

Martens and Medabalimi~\cite{MartensMedabalimi14} prove that the uniform distribution on spanning trees of the complete graph $K_m$ cannot be represented by decomposable and complete (smooth) sum-product networks of size below $2^{m/30240}$ for $m\ge20$, and the same holds for approximation ``arbitrarily well'' by networks of bounded size (their Theorem~37). Their proof writes the network as a sum of $s^2$ ``weak'' functions and shows that each is nonzero on only an exponentially small fraction of the trees; it is a support argument. It does not give a bound at a fixed total-variation error, and the exponent is the number of vertices, which is the square root of the number of variables. Zhang, Holtzen and Van den Broeck~\cite{ZhangHoltzenVdB20} restate this result (their Theorem~5), observe that the uniform spanning tree of $K_n$ is a determinantal point process (Theorem~4), and prove that some determinantal point processes have no polynomial-size deterministic PC. Zhang, Juba and Van den Broeck~\cite{ZhangJubaVdB21} introduce probabilistic generating circuits, which allow subtraction, show that they subsume determinantal point processes, and cite the two results above for the corresponding separation from decomposable PCs. Agarwal and Bl\"aser~\cite{AgarwalBlaser24} and Broadrick, Zhang and Van den Broeck~\cite{BroadrickZhangVdB24} study the power of such circuits further.

Theorem~\ref{thm:main} instead gives a fixed-error bound without assuming determinism or exact representation. It holds for sparse bounded-degree graphs, with an exponent linear in the number of variables. Section~\ref{sec:limits} explains why our method says nothing about $K_n$.

\subsection{Constant-error lower bounds for probabilistic circuits}

Leland and Choi~\cite{LelandChoi25} prove that a \emph{deterministic} decomposable PC that approximates a certain distribution within total variation $\eps\le\frac1{16}-\Omega(1/\mathrm{poly}(n^2))$ has size $2^{\Omega(n)}$ (their Theorem~4.1). The target is built from Sauerhoff's function~\cite{Sauerhoff03} and is itself represented by a polynomial-size nondeterministic decomposable PC. Their proof prunes inputs of small probability, which uses determinism, and reduces to weak approximation by d-DNNFs. Our lower bound does not assume determinism, and our target has a polynomial-size generating circuit with subtraction rather than a small nonnegative one.

The dimer paper~\cite{Blaschke26Dimers} proves an $\exp(\Theta(\sqrt n))$ constant-error bound for uniform domino tilings of the square grid with $n$ edges, and its Section~10.4 notes a stronger generic separation: the uniform distribution on a suitable random linear code, which has easy inference by Gaussian elimination, needs $\exp(\Omega(n))$ size at constant error, by the rank bounds of de Colnet and Mengel~\cite{DeColnetMengel20}. That example is generic rather than natural, and its tractability comes from linear algebra over $\mathbb F_2$. De Colnet and Mengel study weak and strong approximation in knowledge compilation, where a Boolean function is approximated by another Boolean function. Their results do not give ours as a black-box corollary. A distribution within small total variation of $P$ can have full support (take $(1-\delta)P+\delta U$ with $U$ uniform), so approximating $P$ says nothing about approximating its support as a set. Their robust transfer also uses covers by disjoint rectangles, which is why it targets deterministic circuits, while Lemma~\ref{lem:overlap} allows overlapping terms with different splits.

Squared and subtractive circuits~\cite{LoconteEtAl24,WangVdB25} are another way to go beyond nonnegative PCs; we do not consider them.

\subsection{Knowledge compilation and communication complexity}

Bova, Capelli, Mengel and Slivovsky~\cite[Thm.~6]{BovaCapelliMengelSlivovsky16} show that a DNNF of size $s$ yields a cover of its models by at most $s$ balanced rectangles, and use multi-partition communication complexity~\cite{DurisEtAl04} to prove DNNF lower bounds. Sauerhoff~\cite{Sauerhoff03} studies approximation of Boolean functions by such rectangles. Corollary~\ref{cor:support} is a rectangle-cover argument of this kind, with the rectangle bound coming from Corollary~\ref{cor:rectangle}.

The information-complexity method of Bar-Yossef, Jayram, Kumar and Sivakumar~\cite{BJKS04} uses Hellinger distance, the Bhattacharyya coefficient for products of distributions, and the cut-and-paste property of protocols. See also later work such as Jayram's~\cite{Jayram09}. Braun and Pokutta~\cite{BraunPokutta16} use Hellinger-type arguments for robust nonnegative-rank bounds for unique disjointness. Our chain rule (Lemma~\ref{lem:chain}) is a case of Exercise~I.42(b) of Polyanskiy and Wu~\cite{PolyanskiyWu25}, and the identification of the optimized affinity with a matrix norm is Lemma~6 of Lapidoth and Pfister~\cite{LapidothPfister19}. The forbidden outcome $u_i=v_i=1$ of our cycle bits has the form of set disjointness, which connects the argument to the Hellinger methods for disjointness and AND~\cite{BJKS04,Jayram09}. We did not find the combination with a spatial Markov property that we use here, but given how short the argument is, we would not be surprised if similar arguments appear elsewhere.

\subsection{Monotone arithmetic circuits}

Jerrum and Snir~\cite{JerrumSnir82} proved $2^{\Omega(n)}$ monotone lower bounds for the directed spanning-tree polynomial of the complete graph, which has $\Theta(n^2)$ variables. Chattopadhyay, Datta, Ghosal and Mukhopadhyay~\cite[Thm.~1.1]{CDGM22} prove that for a sufficiently large constant $d$ and every $d$-regular expander with second eigenvalue at most $d^{1-\epsilon}$, every monotone circuit for the directed spanning-tree polynomial (one variable per arc, trees directed towards a root) has size $2^{\Omega(n)}$. This is the closest exact analogue of our result that we found. It differs in three ways: it is exact rather than approximate; it concerns the directed polynomial, from which the undirected one is obtained by identifying the two arc variables of each edge (so lower bounds do not transfer from directed to undirected); and it needs large degree, while our graphs have any fixed degree but must be linked. Related monotone lower bounds via communication complexity are in~\cite{ChattopadhyayDattaMukhopadhyay21}. Since the undirected polynomial is a substitution instance of the directed one, a monotone lower bound for the undirected polynomial implies one for the directed polynomial of the same graph, but not conversely. So Corollary~\ref{cor:support}(b) also gives directed bounds for our families. For the undirected polynomial, Jukna and Seiwert~\cite[Thm.~2]{JuknaSeiwert19} prove a $2^{\Omega(\sqrt n)}$ monotone arithmetic lower bound on the complete graph $K_n$, and the exact uniform-spanning-tree lower bound of Martens and Medabalimi~\cite[Thm.~37]{MartensMedabalimi14} also transfers to monotone circuits for the complete-graph polynomial, since replacing $z_e$ by the leaf $[x_e=1]$ and padding with $[x_e=0]$ turns such a circuit into a smooth decomposable PC with polynomial overhead. So exponential exact lower bounds for undirected spanning-tree polynomials are not new in general. What Corollary~\ref{cor:support}(b) adds is bounded-degree families with an exponent linear in the number of variables. A targeted search found no earlier result of this form, but we cannot rule one out.

\subsection{Product approximations of uniform spanning trees}

M\'esz\'aros and Peled~\cite{MeszarosPeled26} prove that for a graph of minimum degree $k\ge2$ there is a coupling of the uniform spanning tree with the random $1$-out subgraph (each vertex picks one incident edge independently) whose expected number of differing edges is $O(n\log(k)/\sqrt k)$. They prove analogous couplings for a wide class of determinantal processes. This does not conflict with our result, though the $1$-out subgraph is close to a product distribution. Their bound is on the expected number of edits, a Wasserstein-type distance, not on total variation; a small expected edit distance per vertex is compatible with total variation close to one. The $1$-out law is a product over vertices, not over a split of the edges, since two vertices can choose the same edge. And at fixed degree the bound is a fixed fraction of $n$, which is small only when $k$ is large. Our result is for fixed degree and total variation.

\subsection{Expanders}

For background on expanders and explicit constructions see Hoory, Linial and Wigderson~\cite{HLW06} and Lubotzky, Phillips and Sarnak~\cite{LPS88}. Our geometric lemma (Lemma~\ref{lem:mixed}) is a giant-component argument of a standard kind.
